\documentclass[11pt]{article}
\usepackage[T1]{fontenc}
\usepackage{textcomp}
\usepackage[margin=0.75in]{geometry}
\usepackage{amsmath,amssymb,amsthm}
\usepackage{array,booktabs}
\usepackage{hyperref}
\setlength{\parindent}{0pt}
\newtheorem{theorem}{Theorem}
\theoremstyle{definition}
\newtheorem{definition}{Definition}
\title{Flow Proof Artifact: prop-21-internal-lines-less-than-sides}
\author{Generated by \texttt{flow doc proof}}
\date{Flow Proof Book}
\begin{document}
\maketitle
If two lines are drawn from the ends of a side to an interior point, their sum is less than the sum of the other two sides.\\[0.5em]
\textbf{Source.} euclid — Elements, Book I, Proposition 21\\[0.5em]
\bigskip
\noindent{\large \textbf{Derived fact 1.} \textit{If two lines are drawn from the ends of a side to an interior point, their sum is less than the sum of the other two sides} \hfill the Euclidean plane \cdot Euclid Book I \cdot Proposition 21: two internal lines from the base are shorter than the other sides}\par
\smallskip\noindent\textit{Coordinate: the Euclidean plane \cdot Euclid Book I \cdot Proposition 21: two internal lines from the base are shorter than the other sides}\par
\smallskip\noindent\textit{Source: euclid — Elements, Book I, Proposition 21}\par
\smallskip\noindent\textit{Built on: proposition 20: the sum of two sides exceeds the third, for Book I of the Elements in on the Euclidean plane}\par
\medskip
\noindent\textbf{Goal.} If two lines are drawn from the ends of a side to an interior point, their sum is less than the sum of the other two sides\par
\noindent$\displaystyle \text{BD plus DC less than AB plus AC}$\par
\medskip
\noindent
\begin{tabular}{@{} >{\raggedright\arraybackslash}p{0.06\textwidth} >{\raggedright\arraybackslash}p{0.47\textwidth} >{\raggedleft\arraybackslash}p{0.41\textwidth} @{}}
\textbf{\#} & \textbf{Proof} & \textbf{Mathematics} \\
\midrule
\textbf{1.} & We prove that proposition 21: two internal lines from the base are shorter than the other sides for Book I of the Elements in the Euclidean plane. & \\[0.45em]
\textbf{2.} & Let D = interior\_point\_inside\_triangle\_ABC. & \\[0.45em]
\textbf{3.} & We invoke the derived fact: triangle\_ABC\_with\_interior\_point\_D. & \\[0.45em]
\textbf{4.} & We invoke the derived fact governing Book I of the Elements in the Euclidean plane: proposition 20: the sum of two sides exceeds the third, for Book I of the Elements in on the Euclidean plane. & \\[0.45em]
\textbf{5.} & From step 2, step 3, and step 4, we can deduce that AB plus AD greater than BD. & $\displaystyle \text{AB plus AD greater than BD}$ \\[0.45em]
\textbf{6.} & We can deduce that AC plus AD greater than CD. & $\displaystyle \text{AC plus AD greater than CD}$ \\[0.45em]
\textbf{7.} & We can deduce that BD plus DC less than AB plus AC. Hence proven. & $\displaystyle \text{BD plus DC less than AB plus AC}$ \\[0.45em]
\bottomrule
\end{tabular}
\label{thm:Geometry-Euclid Book I-proposition_21_two_internal_lines_from_the_base_are_shorter_than_the_other_sides}
\par\medskip\hrule
\end{document}
